\section{Discussion}
\textit{Interpreting transfer asymmetry.} The FIT partitions show descriptive source differences: HP benign URLs are short (median \res{medLenBBenign} characters) and mostly HTTPS (\res{httpsBBenign}\%), whereas HP phishing URLs are longer (\res{medLenBPhish}) with more dots and subdomains (medians \res{medDotsBPhish} and \res{medSubBPhish}). In EBBU, phishing URLs have medians of \res{medSubAPhish} subdomains and \res{medDotsAPhish} dots and are rarely HTTPS (\res{httpsAPhish}\%), but neither are its benign URLs (\res{httpsABenign}\% HTTPS). These differences are consistent with collection-specific cues contributing to the observed asymmetry; they do not establish which feature caused an error. Rashid~\textit{et al.}~\cite{rashid2024} already discuss feature-distribution differences, so identifying source shift is not our novelty. The two datasets differ in collection process, period and prevalence; our transfer drop describes this benchmark pair rather than isolating a causal domain-shift effect.

\textit{What does stability buy?} STABLE-PI prioritises features with more positive top-$q$ fold votes; it does not require every selected feature to be stable in all folds and does not guarantee invariance across datasets. It outperformed mean ranking at the same small budget in HP$\rightarrow$EBBU, without a clear advantage in the other direction. The matched-$k$ design separates this ranking comparison from the number of retained features. ALL nevertheless remained the stronger cross-dataset reference. Our empirical insight is therefore conditional: source-domain stability can mitigate the loss from a compact subset in this benchmark, but it does not establish that source-selected compression preserves transfer or that the rule is generally safer.

\textit{What the two-paper extension establishes.} We connect the classifier/representativeness questions in~\cite{sanchez2022} with the transfer problem in~\cite{rashid2024} through a controlled source-only evaluation. Both primary papers satisfy the five-year recency limit; original EBBU provenance is retained in the repository. The contribution lies in the auditable protocol and the ranking--budget--cost comparison. We have not run either original system under identical conditions, so we cannot attribute improvements over their published results to our approach. The two primary papers satisfy the revised project count and recency conditions; they do not establish universal novelty, and relevant adjacent work remains acknowledged.

\textit{Uncertainty.} The bootstrap intervals resample target domains for fixed fitted models; they do not include resampling or split-induced training variability. Table~\ref{tab:contrasts} shows that per-seed differences vary, so the size of the advantage is less stable than its observed sign. Three estimator seeds on a single split do not establish a population-level training effect, and we make no family-wise significance claim. Additional independently collected holdout data and repeated training splits would strengthen robustness claims, but must be planned under a new protocol before opening a new holdout.

\textit{Cost.} Lexical models are already cheap; reducing 37 to 8--16 features trims extraction work but the shared parsing step and the per-call overhead of the model remain, so savings are modest relative to the accuracy given up on the hard direction. ``Lightweight'' is therefore relative to the ALL model on the reported hardware.

\textit{Threats to validity.} (i) Both datasets are older public collections; results do not establish contemporary or zero-day robustness. (ii) The symmetric purge removes shared popular domains and changes the evaluation population. (iii) Both selectors share an RF proxy, which may favour RF. (iv) Correlated lexical features can depress PI. (v) Hyper-parameters were fixed, not optimised. (vi) Numeric HP label encoding was verified indirectly through the repository's copy of EBBU; complete primary-paper checks remain pending. (vii) The freeze is a procedural control, signed by an agent on the author's instruction. (viii) Some saved SHAP outputs fail the additivity tolerance, so attributions cannot substantiate causal or fully validated explanation claims. The notebook audit also identified direct TEST-vault access in the frozen legacy SHAP analysis; the repaired notebook does not execute that path, but a corrected analysis must follow the versioning rules.
